% !TeX spellcheck = en_US
\documentclass[french]{yLectureNote}

\title{Électromagnétisme 3}
\subtitle{Physique}
\author{Paulhenry Saux}
\date{\today}
\yLanguage{Français}
%jesse.groenen@cemes.fr
\professor{M.Brut}
\usepackage{graphicx}%----pour mettre des images
\usepackage[utf8]{inputenc}%---encodage
\usepackage{geometry}%---pour modifier les tailles et mettre a4paper
%\usepackage{awesomebox}%---pour les boites d'exercices, de pbq et de croquis ---d\'esactiv\'e pour les TP de PC
\usepackage{tikz}%---pour deiffner + d\'ependance de chemfig
% \usepackage{tabularx}%---pour dimensionner automatiquement les tableaux avec variable X
\usepackage{awesomebox}%---Pour les boites info, danger et autres
\usepackage{menukeys}%---Pour deiffner les touches de Calculatrice
\usepackage{fancyhdr}%---pour les en-t\^ete personnalis\'ees
\usepackage{blindtext}%---pour les liens
\usepackage{hyperref}%---pour les liens (\`a mettre en dernier)
\usepackage{caption}%---pour la francisation de la l\'egende table vers Tableau
\usepackage{pifont}
\usepackage{array}%---pour les tableaux
\usepackage{lipsum}
\usepackage{yFlatTable}
\usepackage{multicol}
\usepackage{tabularx}
\usepackage{booktabs}
\newcommand{\Lim}[1]{\lim\limits_{\substack{#1}}\:}
\renewcommand{\vec}{\overrightarrow}
\newcommand{\N}[0]{\mathbb{N}}
\newcommand{\dd}{\mathrm{d}}
\newcommand{\norm}[1]{||\vec{#1}||}
\newcommand{\fo}{\psi(\vec{r},t)}
\newcommand{\foe}{\psi(\vec{r},t)\*}
\newcommand{\HH}{\hat{H}}
\newcommand{\hb}{\hbar}
\newcommand{\lap}{\nabla^2}
\newcommand{\lapcc}{\frac{\partial^2 }{\partial x^2}+\frac{\partial^2 }{\partial y^2}+\frac{\partial^2 }{\partial z^2}}
\newcommand{\E}{\vec{E}(M)}
\newcommand{\B}{\vec{B}(M)}
\newcommand{\V}{V(M)}
\newcommand{\er}{\vec{e_{\rho}}}
\newcommand{\ep}{\vec{e_{\varphi}}}
\newcommand{\pip}{\pi^+}
\newcommand{\pim}{\pi^-}
\newcommand{\mv}{\mathcal{V}}
\DeclareMathOperator{\Div}{Div}
\newcommand{\rot}{\vec{\mathrm{rot}}}
\newcommand{\grad}{\vec{\mathrm{grad}}}
%\setcounter{chapter}{3}
\begin{document}
%\chapter{Puissance et énergie du champ Em}
\setcounter{chapter}{4}
\chapter{Ondes électromahnétiques dans un milieu matériel}
\section{Équations de Maxwell}
\subsection{Équations de Maxwell dans un milieu matériel}
On retrouve les équations de Maxwell. On explicite simplement $\rho = \rho_{libre}+\rho_P+\rho_{autres}$ et $\vec{j} = \vec{j_{libre}}+\vec{j_P}+\vec{j_M}$

avec 
\begin{itemize}
    \item $\rho_P = -\Div \vec{P}$
    \item $\vec{j_P} = \frac{\partial \vec{P}}{\partial t}$
    \item $\vec{j_M} = \rot \vec{M}$
\end{itemize}
En introduisant $\vec{D}$ et $\vec{H}$, on a dans un milieu matériel
\begin{itemize}
    \item $\rot \vec{H} - \frac{\partial \vec{D}}{\partial t}=\vec{j_{libre}}$
    \item $\rot \vec{E}+\frac{\partial \vec{B}}{\partial t} = \vec{0}$
    \item $\Div \vec{D} = \rho_{libre}+\rho_{autres}$
    \item $\Div \vec{B} = 0$
\end{itemize}
\subsection{Relations de passage}
On utilise un cylindre au travers de la surface et on fait tendre le flux vers 0. 
Le flux de D au travers de la surface latérale vaut donc 0.
On transforme enfin les intégrales en $\Delta S$

%TODO vérfier sur les diapositives que les relations de passage sont bien les bonnes

%Attention, on n'a pas besoin de refaire les démos, pour le cours jusqu'ici diapostive 5.6 pour le programme de révision du CC2
\section{Grandeurs énergétiques}
Dans le vide :
\begin{itemize}
    \item Densité volumique dénergie em : $w = \frac{\epsilon_0 E^2}{2}+ \frac{B^2}{2\mu_0}$
    \item Vecteur de Poyting : $\Pi = \vec{E}\wedge \frac{\vec{B}}{\mu_0}$ défini à partir du champ des réels 
    \item Pour une onde progressive dans la direction u : $\vec{Pi} = wc\vec{u}$ : Pi est un courant volumique d'énergie em.
\end{itemize}

Dans le milieu matériel linéaire :
\begin{itemize}
    \item Densité volumique dénergie em : $w = \frac{\vec{D}\cdot \vec{E}}{2}+\frac{\vec{H}\cdot \vec{B}}{2}$
    \item Vecteur de Poyting : $\Pi = \vec{E}\wedge \vec{H}$ défini à partir du champ des réels 
\end{itemize}
On introduit la puissance moyenne rayonnée à travers une surface S :
$$P = \iint_S <\vec{\Pi}>_{temps}\cdot \vec{\dd S}$$
L'intensité vaut alors $ I = \frac{P}{S}$
\section{Ondes électromagnétiqus dans un milieu Diélectrique}
\subsection{Ondes monochormatiques dans un milieu diélectrique LHI}
On consdidère des milueux localment neutres, isolants et non magnétique :
\begin{itemize}
    \item $\rho_{libre} =0$
    \item $\vec{j_{libre}} = 0$
    \item $\vec{M} = \vec{0}$
    \item $\vec{B} = \mu_0\vec{H}$
\end{itemize}
On introduit la relation constitutive du milieu : $\vec{\underline{D}} = \epsilon_0\underline{\epsilon_r}(\omega)\vec{\underline{E}}$ er $\epsilon_r = 1+ \chi_e(\omega)$

Les équations de Maxwell sont identiques à l'exception que $\epsilon_0\to \epsilon_0\epsilon_r$ :Les ondes em qui peuvent se propager dans le milieu sont semblables à celles de propageant dans le vide.

$$\lap \vec{\underline{E}}+\frac{\omega^2}{c^2}\epsilon_r(\omega)\vec{\underline{E}}=\vec{0}$$
avec la vitesse de propagation : $\frac{c}{\sqrt{\epsilon_r}}$
\subsection{Ondes planes progressives monochromatiques}
On est dans le même type de milieu que dans la section précédente et les ondes ont la même aplitude en tout point d'un même plan, i.e. des ondes planes.

Les équations de maxwell sont de la forme :
\begin{itemize}
    \item $\vec{k}\cdot \vec{E} = \vec{k}\cdot \vec{B} = 0$
    \item $\vec{k}\wedge \vec{E} = \omega \vec{B}$
    \item $\vec{k}\wedge \vec{B} = -\omega \mu_0\epsilon_0\epsilon_r(\omega)\vec{E}$
\end{itemize}
 On peut réécrire l'eq de propagation pour une onde plane :
 $$k^2 \vec{\underline{E}} = \frac{\omega^2}{c^2}\epsilon_r(\omega)\vec{\underline{E}}$$
 
 Donc si 
 \begin{itemize}
    \item $\epsilon_r$ réel et strictement positif : milieu transparent, l'onde se propage sans atténuation de l'amplitude : $k$ réel
    \item $\epsilon_r$ dépend de $\omega$ : milieu dispersif
    \item $\epsilon_r =cst$ : milieu non dispersif, c'est le cas du vide.
 \end{itemize}
On introduit la vitesse de phase de l'onde :
$v_{\varphi}(\omega) = \frac{\omega}{k} = \frac{c}{\sqrt{\epsilon_r(\omega)}}$ : La vitesse de phase est plus petite que $c$.

On peut écrire $n(\omega) = \sqrt{\epsilon_r}$.

Dans un milieu matériel, $\lambda = \frac{2\pi}{k} = \frac{\lambda_0}{n(\omega)}$

Le nombre d'onde vaut alors $ k = \frac{\omega}{c}\sqrt{\epsilon_r(\omega)}$
\subsection{Aspects énergétiques}
$$\vec{\Pi} = n(\omega)c\epsilon_0E^2\vec{u}$$
On utilise ici le champ réel !

% \subsection{Milieu diélectrique LHI non magnétique, non chargé, isolant, transparent et non dispersif}
% \begin{itemize}
%     \item $\omega = \frac{ED}{2}+\frac{B^2}{2\mu_0} = \frac{\epsilon E^2}{2}+\frac{B^2}{2\mu_0} = \epsilon_0n^2E^2$
%     \item $\vec{\pi} = \frac{c}{n}w\vec{u}$
%     \item %TODO finir de noter
% \end{itemize}
\subsection{Milieux absorbants}
On a $\epsilon_r$ complexe donc $\vec{k}$ complexe.

% Eq de maxwell, de propagation et de dispersion dans ce cas.

% On note $\vec{k} = (k'+ik'')\vec{u}$

% On a $\vec{E} = \vec{E_0}e^{-k''z}e^{i(k'z-\omega t)}$


% L'indice optique est aussi complexe : $\underline{n} = n'+in''$

% On a $\epsilon_r' = n'^2-n''^2, \epsilon_r''=2n'\cdot n''$

% On a $\underline{n} = \frac{c}{\omega}\underline{k}$

% On a $$<\vec{\Pi}> = \frac12 n' c\epsilon_0E_0^2e^{-2k''z}\vec{e_z}$$

% Si $\epsilon_r'>0$, il y a propagation sans atténuation, milieu transparent, dans le cas contraire, il n'y a plus de propagation, que de l'atténuation.

% Si $\epsilon_r = 0$

% Dans ce cas, il n'y a pas de composante B, donc l'ondes est purement électrique.

% Comme $\vec{k}\wedge \vec{E} = \vec{B} = 0$, l'onde n'est plus transverse mais longitudinale.
\begin{itemize}
    \item Si $\underline{\epsilon_r} = \epsilon_r'$ : 
\begin{itemize}
    \item Si $\epsilon_r' >0$ : propagation sans atténuation : milieu transparent : $e^{i(kz-\omega t)}$
    \item Si $\epsilon_r'<0$ : pas de propagation, que de l'atténuation : $e^{-kz}e^{-i\omega t}$ 
\end{itemize}
\item Dans le cas contraire,  absorption, propagation et atténuation
\end{itemize}

\subsection{Réflexion et transmission d'une onde à l'interface entre 2 diélectriques}
\subsubsection{Position du problème}
On choisit 2 milieux diélectriques LHI non dispersifs, non absorbants, non magnétiques, non chargés et très bon isolants.7
\subsubsection{Relations de passage}
On utilise les relations de passagepour relier les 2  champs du milieu 1 et 2.
\begin{itemize}
    \item $\vec{n_{12}}\cdot (\vec{D_2}-\vec{D_1})=0$
    \item $\vec{n_{12}}\cdot (\vec{B_2}-\vec{B_1})=0$
     \item $\vec{n_{12}}\wedge (\vec{E_2}-\vec{E_1})=0$
    \item $\vec{n_{12}}\wedge (\vec{H_2}-\vec{H_1})=0$
\end{itemize}
Il y a continuité à l'interface des composantes normales de $\vec{D},\vec{B}$ et des composantes tangentielles de $\vec{D}, \vec{E}$
\subsubsection{Relation de Snell-Descartes}
Les vecteurs do'nde $\vec{k_r}, \vec{k_t}$ sont contenus dans le plan d'incidence déifni par le vecteur d'onde $\vec{k_i}$ et la normale à la surface.
\subsubsection{Formules de Fresnel}
La polarisation est la dirrection du champ électrique

Il faut envisager 2  cas :
\begin{itemize}
    \item Polarisation contenue dans le plan d'incidence
    \item Polarisation perpendiculaire au plan d'incidence
\end{itemize}
\subsection{Schémas Bilan}
\includegraphics[scale=0.5]{resume.pdf}
\end{document}
